% =====================================================================
% GreenStep Challenge Admin -- Project Report
% CISC 480 Senior Capstone, Spring 2026
% =====================================================================
\documentclass[11pt,letterpaper]{article}

\usepackage[margin=1in]{geometry}
\usepackage{graphicx}
\usepackage{float}
\usepackage{hyperref}
\usepackage{xcolor}
\usepackage{titlesec}
\usepackage{longtable}
\usepackage{array}
\usepackage{booktabs}
\usepackage{enumitem}
\usepackage{caption}
\usepackage{tocloft}
\usepackage{lastpage}
\usepackage{fancyhdr}

\hypersetup{
  colorlinks=true,
  linkcolor=blue!60!black,
  urlcolor=blue!60!black,
  citecolor=blue!60!black,
  pdftitle={GreenStep Admin Console -- Project Report},
  pdfauthor={Team 2: Josh Dunlap, Eli Goldberger, Khue Vo, Rudy Vergara}
}

\pagestyle{fancy}
\fancyhf{}
\lhead{\small GreenStep Admin Console -- Project Report}
\rhead{\small CISC 480, Spring 2026}
\cfoot{\thepage{} of \pageref{LastPage}}
\renewcommand{\headrulewidth}{0.4pt}

\captionsetup{font=small,labelfont=bf}

\title{
  \vspace{-1cm}
  \textbf{GreenStep Challenge Admin Console} \\[0.5em]
  \large Project Report
}
\author{
  Team 2 \\
  Josh Dunlap \quad Eli Goldberger \quad Khue Vo \quad Rudy Vergara \\[0.5em]
  CISC 480: Senior Capstone, Section 01 \\
  University of St.\ Thomas \\
  Community Partner: Minnesota Pollution Control Agency (MPCA)
}
\date{May 22, 2026}

\begin{document}
\maketitle
\thispagestyle{empty}

\vfill
\noindent
\textbf{Document Type:} Project Report \\
\textbf{Project:} GreenStep Challenge Admin Console \\
\textbf{Repository:} \url{https://github.com/CISC-480-GreenStep/sp26_team2_greenstepchallenge_admin} \\
\textbf{Live URL:} \url{https://greenstep.auriga.fyi} \\
\textbf{Version:} v0.9.0 (final capstone build) \\
\textbf{Submission Date:} May 22, 2026
\vfill
\newpage

\tableofcontents
\newpage

% =====================================================================
\section{Project Overview}
% =====================================================================

\subsection{Problem Statement}

The Minnesota Pollution Control Agency (MPCA) runs annual sustainability
challenges that encourage staff to log everyday actions such as biking to
work, eliminating single-use plastics, and reducing food waste. Before this
project, the agency's Sustainability Program Coordinator,
Kristin Mroz-Risse, administered each challenge through a sprawling network
of spreadsheets, manual email reminders, and ad hoc reporting. The
spreadsheet workflow scaled poorly across multiple concurrent challenges
(department-level cohorts, agency-wide pushes, Earth Month campaigns), made
it difficult to answer the basic question \emph{``who is actually
engaged?''}, and left no audit trail when participants joined or left a
group.

\subsection{Purpose}

The GreenStep Challenge Admin Console replaces the spreadsheet workflow
with a single web application that lets MPCA staff:

\begin{itemize}[noitemsep]
  \item design challenges and reusable templates,
  \item organize participants into flexible groups (departments, cohorts,
        building-level competitions),
  \item monitor engagement on a customizable dashboard, and
  \item export participation reports for leadership.
\end{itemize}

The application is the back office for the broader GreenStep program; a
sibling mobile application built by Team~1 in the same semester gives the
participants themselves a way to log completed actions against a shared
database.

\subsection{Target Users (User Profile)}

The application exposes three role tiers, enforced both in the React UI
and at the database layer via Row-Level Security:

\begin{itemize}
  \item \textbf{SuperAdmin} -- Kristin and her direct delegates. Full
        control: invite new admins, deactivate or permanently delete users,
        manage every group / challenge / template.
  \item \textbf{Admin} -- department-level coordinators. Can create and
        manage challenges, groups, and templates; cannot create or delete
        top-level users.
  \item \textbf{GeneralUser} -- the rest of MPCA staff. Read-only inside
        the admin app; participation itself happens in the sibling mobile
        application.
\end{itemize}

\subsection{Constraints}

\begin{itemize}[noitemsep]
  \item \textbf{Project duration:} ten weeks of active development by a
        four-person student team.
  \item \textbf{Hosting budget:} zero. The deployment had to fit on free
        tiers (Vercel hobby, Supabase free).
  \item \textbf{Network environment:} MPCA staff inboxes are protected by
        corporate link-scanning, which forced the team to abandon
        magic-link authentication in favor of an 8-digit one-time code.
  \item \textbf{Shared database:} the schema is co-owned with Team~1,
        requiring coordination on table and column shape.
  \item \textbf{Browser support:} modern evergreen browsers only; no
        Internet Explorer 11 support.
\end{itemize}

\subsection{Proposed Solution}

A single-page React 19 application hosted on Vercel, talking directly to
a managed Supabase project (Postgres + Auth + REST + Row-Level Security).
The design intentionally has no custom backend tier; authorization is
enforced inside the database through RLS policies that read the role
claim from the caller's JWT. Two narrow Vercel serverless functions
exist only for operations the anon key cannot perform (inviting a new
user, permanently deleting a user account).

This shape gave the team three benefits relative to the original plan of
a separate FastAPI service:

\begin{enumerate}[noitemsep]
  \item one fewer container to deploy, monitor, and secure;
  \item authorization that cannot be bypassed from the browser console;
  \item the shared database becomes a literal contract between Team~2
        (the admin app) and Team~1 (the mobile app).
\end{enumerate}

\subsection{Technology Stack}

\begin{longtable}{p{0.27\linewidth} p{0.68\linewidth}}
\toprule
\textbf{Layer} & \textbf{Technology} \\
\midrule
Frontend         & React 19 (JavaScript) \\
Build tool       & Vite \\
UI library       & Material UI (MUI) v7 with a custom dark theme \\
Charts           & Recharts \\
Dashboard grid   & \texttt{react-grid-layout} v2 (drag-and-drop) \\
Routing          & React Router v7 \\
Database         & Postgres 15 (via Supabase) \\
Auth             & Supabase Auth (8-digit OTP code over email) \\
Authorization    & Postgres Row-Level Security with role-based policies \\
Privileged ops   & Vercel Serverless Functions (Node.js) with Supabase service-role key \\
Hosting          & Vercel (SPA + functions); Supabase (database + auth) \\
Quality gates    & ESLint flat config, Prettier, file-size cap (300 lines / file) \\
\bottomrule
\end{longtable}

\subsection{Expected Benefits}

\begin{itemize}[noitemsep]
  \item \textbf{For Kristin and MPCA:} a single tool that replaces
        the spreadsheet workflow; cross-challenge reports in seconds;
        clear answer to ``who is engaged?''
  \item \textbf{For challenge participants (indirectly):} cleaner
        program operations, fewer broken emails, faster reminders.
  \item \textbf{For the next student team:} a documented codebase with
        every file under 300 lines, a consolidated baseline schema,
        and the C4 architecture, ER diagram, and sequence diagrams
        already on paper.
\end{itemize}

\newpage

% =====================================================================
\section{Current Functionality \& Features}
% =====================================================================

The capstone build (v0.9.0 plus the May 22, 2026 follow-on fixes
merged in PRs \#79--\#87) ships the following features. Every item
in this list is implemented on the live deploy at
\url{https://greenstep.auriga.fyi}.

\subsection{Authentication and Authorization}

\begin{itemize}[noitemsep]
  \item Passwordless sign-in using an 8-digit one-time code emailed
        through Supabase Auth. Pre-flight check rejects unregistered
        or deactivated emails before sending the code.
  \item \texttt{RequireAuth} route guard with optional \texttt{minRole}
        prop; loading spinner during session restore prevents the
        ``flashed login page'' effect on hard refresh.
  \item Row-Level Security enforces the same role checks at the
        database layer, so even a malicious caller using the anon
        key directly cannot mutate admin data.
  \item Dev shortcut: a ``Quick Login as Kristin (SuperAdmin)''
        button on the login page uses \texttt{signInWithPassword}
        against a seeded dev account. Will be removed before any
        real production hand-off.
\end{itemize}

\subsection{User Management}

\begin{itemize}[noitemsep]
  \item Search and filter the user list by role and group; CSV export.
  \item Per-user detail page with activity log, participation
        history, total points, and per-challenge progress bars.
  \item Inline \emph{Change Group} dropdown for Admins / SuperAdmins.
  \item Activate / Deactivate flow with a \texttt{ConfirmDialog};
        deactivation soft-removes by flipping a status flag and
        revoking the auth credential, preserving the audit trail.
  \item SuperAdmin-only \emph{Permanent Delete} (issue \#67) wired
        through a Vercel API route using the Supabase service-role
        key; cascades \texttt{participation} and \texttt{activity\_logs}
        rows while preserving anything the user authored
        (\texttt{createdBy} becomes NULL).
\end{itemize}

\subsection{Group Management}

\begin{itemize}[noitemsep]
  \item Full CRUD for flexible groups (cities, departments, cohorts).
  \item Group detail page with member list, member add/remove,
        challenges scoped to the group, and bidirectional links to
        both.
  \item Member and challenge counts on the group list are clickable
        and navigate to a filtered view.
\end{itemize}

\subsection{Challenge Management}

\begin{itemize}[noitemsep]
  \item List view with search, status filter, group filter, and CSV
        export.
  \item Create and edit form with a four-color theme picker
        (header background, header text, body background, body text)
        and a live mobile preview rendered alongside the form on
        \texttt{md+} screens.
  \item \emph{Quick Start} panel that lets the admin seed a new
        challenge from either a saved \textbf{Template} or a
        previously \textbf{Archived Challenge} (PR \#72).
  \item Challenge detail view: \texttt{Participants} table, full
        participation log, clickable group link, archive button,
        delete button.
  \item Action CRUD per challenge: add, edit, and delete the
        actions associated with a single challenge (name,
        description, category, points).
\end{itemize}

\subsection{Challenge Templates}

\begin{itemize}[noitemsep]
  \item Reusable challenge templates with pre-configured actions.
  \item Three seed templates: H2O Hero Week, Power Down Challenge,
        Sustainable Commute Week.
  \item Applying a template when creating a new challenge pre-fills
        form fields and stages template actions to attach on
        submit.
\end{itemize}

\subsection{Customizable Dashboard}

\begin{itemize}[noitemsep]
  \item 23 available widgets organized into four categories:
        \emph{Overview Stats}, \emph{Charts \& Visualizations},
        \emph{Tables \& Lists}, and \emph{Reports \& Analysis}.
  \item Drag-and-drop reorder and corner-resize via
        \texttt{react-grid-layout}, gated behind an explicit
        \emph{Customize} edit mode.
  \item Four \emph{Quick Layout Presets}: Default, Executive
        Summary, Analytics Deep Dive, Compact Overview.
  \item Layout persists per user to \texttt{localStorage} and
        survives page refreshes.
  \item HCI-informed affordances: grip handles and dashed borders
        only appear in edit mode (progressive disclosure); a
        snackbar confirms saves (feedback); the user cannot remove
        every widget (error prevention).
\end{itemize}

\subsection{Reports and Comparison Mode}

\begin{itemize}[noitemsep]
  \item Standalone \texttt{/reports} page with filter by challenge
        and date range, a category-breakdown chart, the full
        participation table, and CSV export.
  \item New \emph{Participation Report} widget on the dashboard
        (PR \#68) that mirrors the same filterable table + CSV
        export inside a draggable, resizable widget. Each widget
        instance keeps its own filter state so multiple copies can
        be dropped on the grid for side-by-side comparisons.
  \item \emph{Comparison Mode} (PR \#72) promoted from a nested
        dashboard menu to a dedicated toolbar button, with an
        integrated challenge picker.
\end{itemize}

\subsection{Documentation}

\begin{itemize}[noitemsep]
  \item Top-level \texttt{README.md} (project overview, setup,
        version history) and \texttt{CODING\_GUIDELINES.md}
        (Airbnb-style + project-specific conventions).
  \item \texttt{docs/ARCHITECTURE.md} -- the contributor's map.
  \item \texttt{docs/product-documentation/PRODUCT\_DOCUMENTATION.pdf}
        -- the capstone Product Documentation deliverable: PRD,
        UX/UI spec, Technical Design, Source Code Documentation,
        and QA Activities Document, all in one 24-page LaTeX-built PDF.
  \item \texttt{retrospectives.md} -- three documented
        retrospectives (4Ls, Sailboat, Mad/Sad/Glad).
  \item \texttt{supabase/migrations/} -- consolidated baseline:
        \texttt{001\_initial\_schema.sql} (tables, indexes, RPCs,
        RLS, permissions) and \texttt{002\_seed\_data.sql}.
\end{itemize}

\newpage

% =====================================================================
\section{To-Do List \& Future Enhancements}
% =====================================================================

The team intentionally hands this section off to the spring 2027
student team. Items are grouped by category and tagged with the
GitHub issue or PR that tracks them where one exists.

\subsection{Technical}

\begin{itemize}
  \item \textbf{Drop legacy \texttt{presets} / \texttt{preset\_actions}
        tables} (issue \#88). The rename from ``Preset'' to
        ``Template'' was completed at the UI and API level in
        v0.9.0, but the legacy tables were never dropped. One
        line of code in \texttt{data/api/templates.js:92}
        (\texttt{deleteTemplate}) still queries the legacy table
        and silently no-ops.
  \item \textbf{Reconcile \texttt{challenges.theme} (legacy single
        color) with the four \texttt{bgColor*}/\texttt{txColor*}
        columns added in migration~014} (issue \#89). Dashboard
        widgets read the legacy column; the form writes only the
        four split columns. Either migrate widgets or back-fill
        the legacy column on insert.
  \item \textbf{Sync the dashboard layout to the server} so it
        follows the user across devices, replacing the current
        \texttt{localStorage}-only persistence.
  \item \textbf{Tighten the action $\rightarrow$ category foreign
        key.} The \texttt{actions.category} and
        \texttt{challenges.category} columns store the category
        name as text rather than a foreign key to the new
        \texttt{categories} table. FK normalization is a clean
        future migration.
\end{itemize}

\subsection{Functional}

\begin{itemize}
  \item \textbf{Finish issue \#40 (Combine Dashboard and Report Engine).}
        Slice~1 (the Participation Report widget) landed in PR \#68.
        Slices~2 and~3 -- renaming the sidebar entry from
        ``Dashboard'' to ``Report and Analysis,'' adding a
        \emph{Reporting} quick-layout preset, and retiring the
        standalone \texttt{/reports} route -- are scoped but not
        started.
  \item \textbf{Implement Badges and Tickets pages} (issue \#90).
        Draft pages exist on the closed PR \#80 branch but the
        underlying \texttt{badges} and \texttt{tickets} tables
        return 404 on the live database. Both schema and UI
        finishing work are needed; the PR also stubbed out a
        leaderboard query that must be reverted.
  \item \textbf{Photo moderation flow.} Participation rows already
        have a \texttt{photoUrl} column; the actual upload, review,
        and moderation UI was out of scope for this semester.
  \item \textbf{Push notifications} for challenge reminders. The
        mobile app (Team~1) is the natural surface; the admin app
        would only need to enqueue the messages.
  \item \textbf{Saved report templates.} The standalone Reports
        page filters reset on navigate-away. Persisting named
        filters would let admins jump straight to their saved
        cross-challenge slice.
\end{itemize}

\subsection{Usability}

\begin{itemize}
  \item \textbf{Keyboard-accessible dashboard reorder.} Mouse drag
        works; keyboard-only users currently fall back to the
        Quick Layout Presets.
  \item \textbf{Theme presets for the challenge color picker.}
        The four color pickers default to a fixed palette; an
        admin-curated swatch of MPCA brand colors would speed
        up authoring.
  \item \textbf{Inline validation on the challenge form} for the
        end-date-must-be-after-start-date constraint. Currently
        the database constraint is the source of truth; the form
        surfaces the resulting error message but does not
        pre-validate.
\end{itemize}

\subsection{Future Expansion}

\begin{itemize}
  \item \textbf{Full integration with the Team~1 mobile app.} The
        shared schema is in place but the two apps have not been
        exercised together end-to-end at any scale.
  \item \textbf{Multi-tenant support.} If MPCA's program is ever
        adopted by other state agencies, the schema would need a
        tenant identifier on every table or a separate Supabase
        project per tenant.
  \item \textbf{Advanced analytics.} Trend lines over multiple
        years, cohort comparisons, predictive engagement
        indicators -- all out of scope for the spring 2026
        capstone but natural follow-ons once two or three
        years of participation data accumulate.
  \item \textbf{Internationalization.} The current UI is
        English-only.
\end{itemize}

\subsection{Known Bugs and Limitations}

\begin{itemize}
  \item Repository migration files were drifted from the live
        deployed schema for most of the semester. Reconciled and
        squashed to a two-file baseline in PRs \#85 and \#86.
        Documented so the next team doesn't fall into the same trap.
  \item One template apply path still references the legacy
        \texttt{theme} text field rather than the four split
        colors; folded into issue \#89.
\end{itemize}

\newpage

% =====================================================================
\section*{References}
\addcontentsline{toc}{section}{References}
% =====================================================================

\begin{enumerate}[label={[\arabic*]}]
  \item Minnesota Pollution Control Agency, ``GreenStep Cities
        Sustainability Program.'' [Online]. Available:
        \url{https://www.mnpca.state.mn.us/}. [Accessed: May 21, 2026].
  \item S. Brown, \emph{The C4 model for visualising software
        architecture}. [Online]. Available:
        \url{https://c4model.com/}. [Accessed: May 21, 2026].
  \item Supabase, ``Row Level Security.'' [Online]. Available:
        \url{https://supabase.com/docs/guides/auth/row-level-security}.
        [Accessed: May 21, 2026].
  \item Supabase, ``Database Migrations.'' [Online]. Available:
        \url{https://supabase.com/docs/guides/deployment/database-migrations}.
        [Accessed: May 21, 2026].
  \item Material UI, ``Material UI v7 Documentation.'' [Online].
        Available: \url{https://mui.com/material-ui/}. [Accessed:
        May 21, 2026].
  \item Vercel, ``Functions Overview.'' [Online]. Available:
        \url{https://vercel.com/docs/functions}. [Accessed: May
        21, 2026].
  \item J. Patton, \emph{User Story Mapping: Discover the Whole
        Story, Build the Right Product}. O'Reilly Media, 2014.
  \item W3C, ``Web Content Accessibility Guidelines (WCAG) 2.1.''
        [Online]. Available: \url{https://www.w3.org/TR/WCAG21/}.
        [Accessed: May 21, 2026].
\end{enumerate}

\end{document}
